\begin{lemma}[The prime in the divisor conjecture for a surface]
\label{lem:B-surface-Tate-prime}
\status{Proved}
Let \(Y\) be a smooth projective geometrically integral surface over a
finite field \(k\) of characteristic \(p\).  Surjectivity of
\[
 \operatorname{Pic}(Y)\otimes\Q_\ell\longrightarrow
 H^2(Y_{\bar k},\Q_\ell(1))^{\Gamma_k}
\]
for one \(\ell\ne p\) is equivalent to its surjectivity for every
\(\ell\ne p\).  It is also equivalent to the multiplicity of the
eigenvalue \(1\) of Frobenius on \(H^2(Y_{\bar k},\Q_\ell(1))\)
being \(\operatorname{rank}\operatorname{NS}(Y_{\bar k})^{\Gamma_k}\).
\end{lemma}

\begin{proof}
Put \(N=\operatorname{NS}(Y_{\bar k})\otimes\Q\).  Its intersection
pairing is nondegenerate by the Hodge index theorem.  The Galois action
on \(N\) factors through a finite group and preserves that pairing.
Averaging over that group shows that the restriction to \(N^{\Gamma_k}\)
is nondegenerate: an invariant vector orthogonal to all invariant
vectors is orthogonal to every vector.  Norms of divisors show that
the image of \(\operatorname{Pic}(Y)\otimes\Q\) is precisely
\(N^{\Gamma_k}\).  Divisor cycle classes embed this space into the
Frobenius invariants for every \(\ell\).

Suppose that surjectivity holds for one \(\ell\), and write \(V\) for
the corresponding degree-two twisted cohomology and \(F\) for
Frobenius.  The subspace \(V^F=N^{\Gamma_k}\otimes\Q_\ell\) is
nondegenerate for the cup product.  Its orthogonal complement is
\(F\)-stable and contains no fixed vector.  Thus \(F-1\) is invertible
on that complement; in particular the generalized \(1\)-eigenspace
equals \(V^F\).  The characteristic polynomial of Frobenius on
degree-two cohomology is independent of \(\ell\).  For every other
prime, the algebraic invariant subspace therefore already has the
full dimension of the generalized \(1\)-eigenspace.  This proves the
assertions, including the converse multiplicity criterion.  No
general Frobenius semisimplicity conjecture is used.
\end{proof}

\begin{theorem}[A smaller surface test class: vanishing coherent irregularity]
\label{thm:B-irregularity-zero-surface-gate}
\status{Proved equivalence; the restricted Tate conjecture remains unresolved}
Fix a finite field \(k\) and a prime \(\ell\ne\operatorname{char}k\).
The divisor Tate conjecture at \(\ell\) for every smooth projective
geometrically integral surface over \(k\) is equivalent to the same
assertion restricted to surfaces \(Y/k\) satisfying
\[
                         H^1(Y,\mathcal O_Y)=0.
\]
Consequently every assertion of
Theorem~\ref{thm:B-relative-surface-gate}, and the equivalent universal
relative Brauer assertions proved above, can be tested using only
these surfaces over the finite extensions of \(\F_p\).
This is a proper subclass restriction: for an elliptic curve \(E/k\),
the surface \(E\times\mathbb P^1_k\) has
\(\dim_kH^1(\mathcal O)=1\) and is excluded.  The restricted hypothesis
imposes only the displayed intrinsic vanishing; it does not retain
an arbitrary original surface as auxiliary data.
\end{theorem}

\begin{proof}
The implication to the restricted class is immediate.  For the
converse, Lemma~\ref{lem:B-surface-Tate-prime} lets us use the
prime-independent notation \(T^1\) of
\citet[Conjecture~1.5 and Theorem~1.6, pp.~1967--1968]{DYuan2021Hodge}.
We give the reduction, keeping its base field and its zero-dimensional
case explicit.  The full published proof is in that paper's
Section~4.2, pp.~1996--1997; its preparatory results are
Lemmas~4.2 and~4.4 and Theorems~4.3 and~4.5,
pp.~1994--1996.

Let \(X/k\) be any smooth projective geometrically integral surface.
The finite-field Lefschetz-pencil construction used by Yuan produces
a smooth blowup \(X'\) and a semistable fibration
\(X'\to\mathbb P^1_k\), over this same \(k\).
Write \(J/k(t)\) for the Jacobian of its smooth generic fibre.
Blowups add only algebraic divisor classes; the Artin--Tate
comparison then gives
\[
                  T^1(X)\ \Longleftrightarrow\ T^1(X')
                         \ \Longleftrightarrow\ \operatorname{BSD}(J),
\]
where \(\operatorname{BSD}\) here means equality of the analytic and
Mordell--Weil ranks, not the conjectural leading-coefficient formula.

By Yuan's Theorem~1.1, p.~1965, there is an isogeny over \(k(t)\)
\[
                         J\sim A\times C_{k(t)},
\]
where \(C/k\) is an abelian variety and the Hodge bundle of \(A\)
on \(\mathbb P^1_k\) is ample.  Semi-abelian reduction is preserved
under isogeny, so \(A\) has it everywhere.  The constant factor has
rank zero: a map \(\mathbb P^1_k\to C\) is constant, hence its
Mordell--Weil group is \(C(k)\).  Its analytic rank is also zero,
as the Frobenius eigenvalues of \(C\) have absolute value
\(\sqrt{\#k}\).  Isogeny invariance and additivity of the ranks,
together with the known nonnegativity of each analytic-rank defect,
give \(\operatorname{BSD}(J)\Longleftrightarrow
\operatorname{BSD}(A)\).  If \(A=0\), this already proves
\(T^1(X)\), and nothing further is required.

Assume \(\dim A\ge1\).  Yuan's Theorem~4.5 supplies a projective flat
regular model \(\psi:P\to\mathbb P^1_k\) of \(A^\vee\) with
\[
                    R^1\psi_*\mathcal O_P
                    \simeq\operatorname{Lie}(\mathcal A/\mathbb P^1_k),
\]
where \(\mathcal A\) is the N\'eron model of \(A\).  This bundle is
the dual of its ample Hodge bundle.  It has no global sections.
Since \(\psi_*\mathcal O_P=\mathcal O_{\mathbb P^1_k}\), the
coherent Leray sequence yields \(H^1(P,\mathcal O_P)=0\).
Regularity over the perfect field \(k\) makes \(P\) smooth.

If \(\dim P=2\), set \(Y=P\).  Otherwise successively take
sufficiently ample smooth hypersurface sections over \(k\) until
obtaining a surface \(Y\subset P\).  At each step choose a smooth
closed fibre over a point of sufficiently large residue degree.
Poonen's Bertini density theorem
\citep[Theorem~1.1, author-version p.~2]{DPoonen2004Bertini}
makes it possible to choose the
section simultaneously smooth on the total space and on that fibre:
the first density is fixed and positive, while the second tends to
one with that residue degree.  Smoothness of the generic fibre
follows.  Serre vanishing and duality give, at each cut,
\[
 H^1(Y,\mathcal O_Y)=0,
 \qquad
 H^1(P_\eta,\mathcal O_{P_\eta})
       \lhook\joinrel\longrightarrow H^1(Y_\eta,\mathcal O_{Y_\eta}),
\]
with the evident interpretation for the successive intermediate
sections.  For the first equality the required vanishing is
\(H^2(P,\mathcal O_P(-Y))=0\); for the injection on a generic
fibre it is \(H^1(P_\eta,\mathcal O_{P_\eta}(-Y_\eta))=0\).
The same ampleness ensures geometric connectedness.

Let \(B\) be the Jacobian of the resulting smooth generic curve
\(Y_\eta\).  Pullback of line bundles gives \(A\to B\), whose map
on tangent spaces is the displayed injection.  Its kernel is
finite, and Poincar\'e reducibility makes \(A\) an isogeny factor of
\(B\).  Thus
\[
 T^1(Y)\ \Longrightarrow\ \operatorname{BSD}(B)
       \ \Longrightarrow\ \operatorname{BSD}(A)
       \ \Longrightarrow\ T^1(X).
\]
Our hypothesis supplies the first assertion because
\(H^1(Y,\mathcal O_Y)=0\).  Every construction was over \(k\).
Applying this fixed-field result for every finite \(k/\F_p\)
handles all the constant fields in
Theorem~\ref{thm:B-relative-surface-gate}; viewing a geometrically
connected surface over its field of constants and using induction
of its Galois representation also handles the non-geometrically
connected surfaces appearing there.
\end{proof}

\begin{remark}[Sources and scope of the reduction]
\label{rem:B-irregularity-zero-source-scope}
\status{Verified source scope}
\citet[Remark~4.13, p.~56]{BErtlKellerQin2025} points to the
all-dimensional reduction and the additional restriction above.
The existing Kahn--Morrow reduction supplies the passage from
arbitrary dimension to surfaces; Yuan's Theorem~1.6 supplies the
further vanishing of coherent irregularity.  Yuan's cited
Nguyen thesis and K\"unnemann--Faltings--Chai inputs are used through
his exact published statements and proof; they are not represented
as independently opened here.  De Jong's unpublished reduction note
has not been obtained.  None of these reductions proves the Tate
conjecture for the remaining irregularity-zero surfaces.
\end{remark}

\begin{corollary}[Canonical surfaces suffice, with arbitrary numerical thresholds]
\label{cor:B-canonical-surface-gate}
\status{Proved strict test-class reduction; the restricted Tate conjecture remains unresolved}
Fix a finite field \(k\), a prime \(\ell\ne\operatorname{char}k\),
an integer \(N\ge1\), and a real number \(\varepsilon>0\).
Divisor Tate at \(\ell\) for every smooth projective geometrically
integral surface over \(k\) is equivalent to its restriction to
surfaces \(Y/k\) with
\[
 H^1(Y,\mathcal O_Y)=0,\qquad \omega_Y\text{ very ample},
 \qquad K_Y^2\ge N,\qquad p_g(Y)\ge N,
\]
and, if desired, the simultaneous extra conditions
\[
             c_2(Y)>0,\qquad
             \left|\frac{K_Y^2}{c_2(Y)}-1\right|<\varepsilon.
\]
Here \(c_2(Y)=\deg c_2(T_Y)\).
Thus the preceding universal Brauer and abelian-torsor gates can
be confined to an intrinsic class of canonically embedded surfaces
of general type.  This does not assert Tate for those surfaces.
\end{corollary}
\begin{proof}
Use the construction in the preceding theorem for an arbitrary
surface \(X/k\).  If its ample-Hodge factor \(A\) is zero,
\(T^1(X)\) is already proved.  Otherwise let
\(P\to\mathbb P^1_k\) be the smooth projective regular model
of \(A^\vee\) constructed there, with \(H^1(P,\mathcal O_P)=0\).
If \(\dim P=2\), replace \(P\) by \(P\times_k\mathbb P^1_k\),
using projection to the original base.  K\"unneth leaves
\(H^1(\mathcal O)\) unchanged, and its smooth generic fibre has
the same \(\operatorname{Pic}^0=A\).  Thus the new ambient
dimension is at least three in every case.

Perform the high-degree smooth cuts of Yuan's proof only until
reaching a smooth projective threefold \(W\) over \(k\).
They preserve \(H^1(W,\mathcal O_W)=0\), a smooth generic
fibre over \(\mathbb P^1_k\), and the injection from the
original generic-fibre \(H^1(\mathcal O)\).
Choose a very ample line bundle \(L\) on \(W\), with
\(H=c_1(L)\).  The final cut \(Y\in|L^d|\) may be taken
for arbitrarily large \(d\), simultaneously smooth on \(W\)
and on one fixed sufficiently high-degree smooth closed fibre,
by the positive-density argument already proved above.
Serre duality and vanishing give
\[
 H^2(W,L^{-d})=0,\qquad
 H^1(W_\eta,L^{-d}|_{W_\eta})=0
 \quad(d\gg1).
\]
Consequently \(H^1(Y,\mathcal O_Y)=0\), and the resulting
generic Jacobian still contains \(A\) up to isogeny.
The same high-degree choices ensure geometric connectedness.

There is one more threshold, compatible with all those above:
\(\omega_W\otimes L^d\) is very ample for every sufficiently
large \(d\)~\cite[Lemma~29.40.8, tag 0FVC]{BStacksCanonicalCuts}.
Taking determinants in the exact conormal sequence gives
\[
                  \omega_Y=(\omega_W\otimes L^d)|_Y,
\]
which is therefore very ample.  This determinant argument is
valid in every characteristic; the conormal sequence is short
exact because both \(W\) and \(Y\) are smooth
\cite[section~37.63, tag 06BB]{BStacksCanonicalCuts}.

All numerical conditions can be imposed in this same cut, not
as separate later reductions.  Put \(K=c_1(\omega_W)\).
Adjunction and the normal-bundle sequence give
\[
 \begin{aligned}
 K_Y^2&=d^3H^3+2d^2KH^2+dK^2H,\\
 c_2(Y)&=d^3H^3+d^2KH^2+d\,c_2(W)H.
 \end{aligned}
\]
Their common leading coefficient is \(H^3>0\), so
\(K_Y^2\to\infty\), \(c_2(Y)>0\) for large \(d\),
and \(K_Y^2/c_2(Y)\to1\).
The exact divisor sequence also gives
\[
 \chi(\mathcal O_Y)
   =\chi(\mathcal O_W)-\chi(L^{-d})
   =\frac{H^3}{6}d^3+O(d^2).
\]
Since \(Y\) is geometrically connected and \(H^1(\mathcal O_Y)=0\),
Serre duality identifies
\(p_g(Y)=\chi(\mathcal O_Y)-1\), which tends to infinity.
Increasing the one threshold for \(d\) therefore gives all
the displayed inequalities.  The proof of the preceding theorem
still gives \(T^1(Y)\Rightarrow T^1(X)\), over the original
fixed field \(k\).  This proves the equivalence.
\end{proof}

This restriction is intrinsic: it records no arbitrary original
surface as a blowup centre or other additional datum.
For example, \(\mathbb P^2_k\) has \(H^1(\mathcal O)=0\)
but is excluded by the canonical ampleness condition.  The
numerical refinements above are consequences of the same
high-degree-cut argument and are counted as part of this one
reduction, not as successive movements of the conjectural gate.
